\documentclass[preprint,3p,number]{elsarticle}

\usepackage{amsmath,amssymb}

\title{Finite Element Discretization of Two-Dimensional Natural Convection}
\author{}
\date{}

\begin{document}

\maketitle

\section{Finite Element Discretisation}

\subsection{Weak Formulation}

Let $\Omega$ denote the computational domain and
$\partial\Omega$ its boundary. The velocity vector is
$\vec{U}=(u,v)$. Let $q$, $\vec{W}$,
and $\theta$ denote the test functions for the continuity,
momentum, and energy equations, respectively. The symbol
$\vec{W}$ denotes the vector momentum test function,
whereas the unadorned $v$ denotes the vertical velocity
component.

\subsubsection{Continuity Equation}

The continuity equation is

\begin{equation}
\nabla\cdot\vec{U}=0.
\end{equation}

Multiplying by the scalar test function $q$ gives

\[
q\,\nabla\cdot\vec{U}=0.
\]

Integrating over the domain gives

\begin{equation}
\boxed{
\int_{\Omega}q\,\nabla\cdot\vec{U}\,d\Omega=0.
}
\end{equation}

This is the weak statement of continuity. Integration by
parts is not required for this particular formulation.
Although the velocity is a vector, the continuity equation
is a scalar equation.

\subsubsection{Momentum Equations}

The two momentum equations can be combined into one
vector equation:

\begin{equation}
\rho_0\left[
(\vec{U}\cdot\nabla)\vec{U}
\right]
=-\nabla p+\mu\nabla^2\vec{U}+\mathbf{f}_b.
\end{equation}

where $\vec{U}=(u,v)$ is the velocity vector,
$p$ is the pressure, $\rho_0$ is the reference density,
$\mu$ is the dynamic viscosity, and $\mathbf{f}_b$
is the buoyancy force per unit volume, given by

\begin{equation}
\mathbf{f}_b=
\begin{pmatrix}
0\\
\rho_0g\beta(T-T_{\mathrm{ref}})
\end{pmatrix}.
\end{equation}

The buoyancy term appears only in the vertical component
under the chosen coordinate system.

These are the strong forms: the original differential
equations before applying the weak-form procedure.

\paragraph{Step 1: Multiply by a vector test function}

Let $\vec{W}$ be a vector test function.
Take the dot product with it and integrate over $\Omega$:

\begin{equation}
\begin{aligned}
&\int_\Omega \rho_0\vec{W}\cdot
\left[
(\vec{U}\cdot\nabla)\vec{U}
\right]d\Omega\\
&=-\int_\Omega\vec{W}\cdot\nabla p\,d\Omega
+\mu\int_\Omega\vec{W}\cdot
\nabla^2\vec{U}\,d\Omega\\
&\quad+\int_\Omega\vec{W}\cdot
\mathbf{f}_b\,d\Omega.
\end{aligned}
\end{equation}

\paragraph{Step 2: Apply integration by parts}

For the pressure term, use the divergence theorem:

\begin{equation}
\begin{aligned}
-\int_\Omega\vec{W}\cdot\nabla p\,d\Omega
={}&\int_\Omega p\,\nabla\cdot\vec{W}\,d\Omega\\
&-\int_{\partial\Omega}p\,
\vec{W}\cdot\mathbf{n}\,d\Gamma.
\end{aligned}
\end{equation}

For the viscous term:

\begin{equation}
\begin{aligned}
\mu\int_\Omega\vec{W}\cdot
\nabla^2\vec{U}\,d\Omega
={}&-\mu\int_\Omega
\nabla\vec{W}:\nabla\vec{U}\,d\Omega\\
&+\mu\int_{\partial\Omega}
\vec{W}\cdot
\frac{\partial\vec{U}}{\partial n}\,d\Gamma.
\end{aligned}
\end{equation}

Here, $\mathbf{n}$ is the outward unit normal to the
boundary, $d\Gamma$ is a boundary element, and the colon
denotes the sum of products of corresponding matrix
components.

\paragraph{Step 3: Write the weak momentum equation}

Substituting these expressions into the integrated
momentum equation and rearranging gives

\begin{equation}
\boxed{
\begin{aligned}
&\int_\Omega \rho_0\vec{W}\cdot
\left[
(\vec{U}\cdot\nabla)\vec{U}
\right]d\Omega\\
&+\mu\int_\Omega
\nabla\vec{W}:\nabla\vec{U}\,d\Omega\\
&-\int_\Omega p\,\nabla\cdot\vec{W}\,d\Omega
-\int_\Omega\vec{W}\cdot
\mathbf{f}_b\,d\Omega=B_m.
\end{aligned}
}
\end{equation}

where the boundary contribution is

\begin{equation}
\begin{aligned}
B_m={}&-\int_{\partial\Omega}p\,
\vec{W}\cdot\mathbf{n}\,d\Gamma\\
&+\mu\int_{\partial\Omega}
\vec{W}\cdot
\frac{\partial\vec{U}}{\partial n}\,d\Gamma.
\end{aligned}
\end{equation}

For a no-slip wall, the velocity is prescribed and the
momentum test function is zero on that wall. The
corresponding boundary terms vanish there. For other
boundary conditions, the boundary terms must be treated
according to the specified physical conditions.

This is the weak form of the momentum equations.

\subsubsection{Energy Equation}

The energy equation is

\begin{equation}
\vec{U}\cdot\nabla T
=\alpha\nabla^2T.
\end{equation}

\paragraph{Step 1: Multiply by a scalar test function}

Multiply both sides by the scalar test function $\theta$:

\begin{equation}
\theta\,\vec{U}\cdot\nabla T
=\alpha\theta\nabla^2T.
\end{equation}

Integrating over the domain gives

\begin{equation}
\begin{aligned}
&\int_\Omega \theta\,\vec{U}\cdot\nabla T\,d\Omega\\
&=\alpha\int_\Omega \theta\nabla^2T\,d\Omega.
\end{aligned}
\end{equation}

\paragraph{Step 2: Apply integration by parts}

Applying integration by parts to the diffusion term gives

\begin{equation}
\begin{aligned}
\alpha\int_\Omega \theta\nabla^2T\,d\Omega
={}&-\alpha\int_\Omega
\nabla\theta\cdot\nabla T\,d\Omega\\
&+\alpha\int_{\partial\Omega}
\theta\frac{\partial T}{\partial n}\,d\Gamma.
\end{aligned}
\end{equation}

\paragraph{Step 3: Write the weak energy equation}

Moving the diffusion term to the left gives

\begin{equation}
\boxed{
\begin{aligned}
&\int_\Omega \theta\,\vec{U}\cdot\nabla T\,d\Omega\\
&+\alpha\int_\Omega
\nabla\theta\cdot\nabla T\,d\Omega=B_T,
\end{aligned}
}
\end{equation}

where

\begin{equation}
B_T=\alpha\int_{\partial\Omega}
\theta\frac{\partial T}{\partial n}\,d\Gamma.
\end{equation}

For prescribed-temperature boundaries, the test function
$\theta$ is zero on those boundaries. For insulated boundaries,
$\partial T/\partial n=0$. Under these conditions, the
corresponding boundary contributions vanish.

This is the weak form of the energy equation.

\subsection{Finite Element Approximation}

The weak forms are discretised by approximating the
continuous fields using finite element basis functions.
For example,

\begin{align}
u_h(x,y)&=\sum_{j=1}^{N_u}U_j\phi_j(x,y),\\
v_h(x,y)&=\sum_{j=1}^{N_v}V_j\psi_j(x,y),\\
p_h(x,y)&=\sum_{j=1}^{N_p}P_j\chi_j(x,y),\\
T_h(x,y)&=\sum_{j=1}^{N_T}T_j\eta_j(x,y).
\end{align}

Here, the basis functions define the spatial approximation,
while the coefficients are the unknown degrees of freedom.

The discrete velocity and pressure fields are substituted
into the weak momentum and continuity equations. Similarly,
the discrete velocity and temperature fields are substituted
into the weak energy equation. The test functions are
chosen from the corresponding finite element spaces.

This procedure produces a finite system of algebraic
equations for the unknown degrees of freedom. Because the model is steady-state, no time-discretisation
is required. The resulting spatially discretised coupled
system is solved directly for the steady velocity, pressure,
and temperature fields.


\end{document}
